\section{Materials and Methods}
\label{sec:method}

Figure~\ref{fig:framework} overviews EARS-MMOEA. A real OpenStreetMap city is cast as an MMOP and
solved by seven cooperating modules; the output is a set of Pareto-equivalent solution
families. Algorithm~\ref{alg:ears} gives the generational loop and Table~\ref{tab:notation}
the notation. This section first presents the core hybrid diversity key, then the supporting
modules and the structured representation used for design tasks.

\begin{figure}[t]
  \centering
  \includegraphics[width=\textwidth]{framework_ai.pdf}
  \caption{EARS-MMOEA framework. A real OpenStreetMap city (left) feeds the generational engine
  (centre), whose core M1 is the within-front diversity selection key $D(i){=}E(i)(1{+}\beta S(i))$;
  the output (right) is a set of Pareto-equivalent station layouts.}
  \label{fig:framework}
\end{figure}

\begin{table}[t]
\centering\small
\caption{Notation.}
\label{tab:notation}
\begin{tabular}{@{}l p{0.66\linewidth}@{}}
\toprule
symbol & meaning \\
\midrule
$N$, $m$ & population size, number of objectives \\
$x_i$, $F(x_i)$ & decision vector and its objective vector \\
$\phi(\cdot)$ & decision-space representation (identity, or a feature map) \\
$\mathrm{SCD}(i)$ & special crowding distance (max of obj.- and dec.-space crowding) \\
$r(i)\in[0,1]$ & niche rarity of $i$ (rarer modes scored higher) \\
$E(i)$ & equivalence diversity $=\mathrm{SCD}(i)\,(1+\rho\,r(i))$, rarity boost $\rho{=}1.0$ \\
$S(i)\in[0,1]$ & normalised $k$NN sparsity of $\phi(x_i)$ among non-dominated solutions ($k{=}3$) \\
$\beta$ & sparsity weight (frozen at $0.5$) \\
$q_o$ & bandit quality of operator $o$; updated by reward $=$ offspring improvement over parent \\
$D(i)$ & hybrid diversity key, Eq.~\eqref{eq:hybrid} \\
\bottomrule
\end{tabular}
\end{table}

\begin{algorithm}[t]
\caption{EARS-MMOEA (one run).}
\label{alg:ears}
\begin{algorithmic}[1]
\Require problem: box bounds $\mathbf{x}^l,\mathbf{x}^u\in\mathbb{R}^{d}$ and evaluator
  $\textsc{Eval}:\mathbb{R}^{N\times d}\!\to\!(\mathbb{R}^{N\times m},\mathbb{R}_{\ge 0}^{N})$
  returning objectives $F$ and constraint violation $\mathrm{CV}$; population size $N\in\mathbb{N}$,
  budget $T\in\mathbb{N}$ (function evaluations); frozen scalars $\beta{=}0.5$, $\rho{=}1.0$,
  $k{=}3\in\mathbb{N}$, re-cluster period $\tau{=}10\in\mathbb{N}$, mode range $[2,5]$,
  cross-mode base $p_0\in[0,1]$, bandit rate $\alpha\in(0,1)$, floor $p_{\min}\in[0,1)$,
  diversity-archive feed probability $p_{\mathrm{feed}}\in[0,1]$
\Ensure $\mathcal{A}_{\mathrm{obj}}$: the non-dominated, equivalence-spread Pareto-set
  approximation (decision vectors in $\mathbb{R}^{d}$ with their objectives)
\State $P \gets N$ uniform samples in $[\mathbf{x}^l,\mathbf{x}^u]$ \Comment{$P\in\mathbb{R}^{N\times d}$}
\State $(F,\mathrm{CV}) \gets \textsc{Eval}(P)$ \Comment{$F\in\mathbb{R}^{N\times m}$, $\mathrm{CV}\in\mathbb{R}_{\ge0}^{N}$}
\State init Pareto archive $\mathcal{A}_{\mathrm{obj}}$, decision-mode archive $\mathcal{A}_{\mathrm{mode}}$,
  diversity archive $\mathcal{A}_{\mathrm{div}}$, operator portfolio $\Pi$ (probability matching),
  $\epsilon$-controller
\State update $\mathcal{A}_{\mathrm{obj}},\mathcal{A}_{\mathrm{mode}},\mathcal{A}_{\mathrm{div}}$ with $P$;\quad
  modes $\gets$ silhouette $k$-means on $\phi(P)$;\quad $g\gets0$
\While{evaluations used $\le T-N$}
  \State $g\gets g+1$;\quad \textbf{if} $g \bmod \tau = 0$ \textbf{then} modes $\gets$ silhouette $k$-means on $\phi(P)$
  \State labels, occupancy entropy $H \gets$ assign $P$ to modes
  \State $p_{\mathrm{cross}} \gets \mathrm{clip}\!\big(p_0 + 0.5(1-H),\,0.05,\,0.9\big)$
  \For{$i = 1$ to $N$} \Comment{reproduction: within-/cross-mode mating + bandit operator}
    \State $(a,b)\gets$ cross-mode pair (two distinct modes) w.p.\ $p_{\mathrm{cross}}$, else within-mode pair
    \State w.p.\ $p_{\mathrm{feed}}$ replace $b$ by a converged solution drawn from $\mathcal{A}_{\mathrm{div}}$
    \State $o_i \gets \Pi.\textsc{Select}()$;\quad $Q_i \gets \text{operator}_{o_i}(P_a,P_b)$
  \EndFor
  \State $(F_Q,\mathrm{CV}_Q)\gets \textsc{Eval}(Q)$
  \For{$i=1$ to $N$} \Comment{credit assignment}
    \State $\textit{reward}\gets \mathbb{1}[\,Q_i\text{ feasible and dominates its base parent}\,]$;\quad
    $\Pi.\textsc{Reward}(o_i,\textit{reward})$
  \EndFor
  \State update $\mathcal{A}_{\mathrm{obj}},\mathcal{A}_{\mathrm{mode}},\mathcal{A}_{\mathrm{div}}$ with $Q$
  \State $R \gets P \cup Q$;\quad advance $\epsilon$-schedule;\quad relax $\mathrm{CV}_R$ if constrained;\quad label $R$ by nearest mode centre
  \State constraint-aware non-dominated sort of $R$
  \State fill $P$ front-by-front; on the splitting front $F_\ell$ keep the highest
  $D(i){=}E(i)(1{+}\beta S(i))$ until $N$ are selected (Eq.~\eqref{eq:hybrid}), in representation $\phi$
\EndWhile
\State \Return $\mathcal{A}_{\mathrm{div}}.\textsc{Report}()$ if enabled, else
  non-dominated$(P\cup\mathcal{A}_{\mathrm{obj}}\cup\mathcal{A}_{\mathrm{mode}})$
\end{algorithmic}
\end{algorithm}

\subsection{Core: hybrid dual-space diversity (M1)}
\label{sec:hybrid}
The environmental selection ranks the splitting non-dominated front by the key
\begin{equation}
  D(i) \;=\; E(i)\,\bigl(1 + \beta\, S(i)\bigr),
  \label{eq:hybrid}
\end{equation}
where $E(i)$ is the \emph{equivalence diversity} (the SCD of solution $i$, the larger of
its objective- and decision-space crowding, boosted by niche rarity so that members of rarer
decision-space modes are scored higher), and $S(i)\in[0,1]$ is the normalised sum of
$i$'s $k$ nearest-neighbour distances in (a problem feature representation of) decision
space (a decision-space sparsity). The weight $\beta$ is frozen at $0.5$ by
validation-only analysis (Section~\ref{sec:param}).

\paragraph{Design rationale: where the diversity term enters} The bonus is always active: with a
finite population $S(i)>0$ almost everywhere, so $D$ does not collapse to $E$. The design works by
\emph{placing} the diversity term, not through a dominance proof.
The key is applied \emph{only after} the constraint-aware non-dominated sort, purely as an
intra-front term that re-orders the splitting front $F_\ell$ (the partial front that overflows the
next population). Because $1+\beta S$ is non-negative and acts only among mutually
non-dominated members of $F_\ell$, no dominated solution is ever promoted above a
non-dominated one; it preserves \emph{front precedence}, re-ordering within a front but never
letting a worse front overtake a better one. This property is shared with NSGA-II-style
front-by-front (intra-front) selection. In the inter-/intra-front
taxonomy~\parencite{javadi2022interfront} our term is an intra-front operation, and our contribution is
the controlled finding that, under a fixed backbone and the \emph{same} unpenalised sparsity
signal, placing it within the front rather than letting it override front precedence (as a global
key would) incurs a smaller convergence loss. CPDEA and MMEA-WI instead let a decision-space signal
influence which solutions count as converged; in our controlled setting this admits
somewhat-unconverged solutions and degrades IGD. The finding is scoped to the unpenalised signal:
a convergence-penalised key behaves differently in the sort (Section~\ref{sec:novelty}), and the
full behaviour of those methods also reflects their distinct operators and archives.
Re-ordering $F_\ell$ can still perturb the search trajectory, so the cost is \emph{small but
not zero}: we measure it. The sparsity bonus alone costs $\le 3\%$ IGD
(the full key A0 vs the no-bonus variant A9, Section~\ref{sec:ablation}); the whole decision-space
stack costs $\approx 7\%$ IGD over an NSGA-II base (the constructive ladder, same section). These
are different comparisons. By ``orthogonal axes'' we mean this selection-time decomposition:
convergence in the rank, diversity in the within-rank factor. The multiplicative form is
interchangeable with an additive one: a controlled
experiment (Section~\ref{sec:novelty}) shows an additive within-front variant
$E+\beta\,\mathrm{med}(E)\,S$ performs identically, while moving $S$ into the dominance sort
significantly degrades convergence. The placement, not the form, carries the benefit; we use
the multiplicative form as the default. The empirical effect of $S$ is quantified by ablation variant A9
(Section~\ref{sec:ablation}): removing it worsens IGDX by $11\%$ (A0 beats A9 $3/5/0$ on IGDX,
Holm-significant on the close-mode instances, a tie elsewhere).

\paragraph{Module formulas} Two distinct weights should not be conflated: the
\emph{rarity boost} $\rho$ (frozen at $1.0$) acts \emph{inside} the equivalence
diversity $E$ (favouring sparse decision-space \emph{modes}), whereas the \emph{sparsity
weight} $\beta$ of Eq.~\eqref{eq:hybrid} multiplies the per-solution decision-space sparsity
$S$ \emph{outside} $E$. For completeness the splitting-front members are scored by
$E(i)=\mathrm{SCD}(i)\,(1+\rho\,r(i))$ with
$r(i)=1-(c_i-1)/c_{\max}\in[0,1]$, where $c_i$ is the occupancy of $i$'s decision-space mode
and $c_{\max}$ the largest occupancy (rarer modes scored higher). $\mathrm{SCD}$ is the
special crowding distance~\parencite{yue2018moring} (the larger of objective- and decision-space
crowding). The sparsity $S(i)$ is the min--max-normalised sum of the $k{=}3$ nearest-neighbour
distances of $\phi(x_i)$ among the front (so $S\in[0,1]$, sparsest $=1$); boundary members
with infinite $\mathrm{SCD}$ keep top priority. The operator bandit (M6) uses
probability matching: quality $q_o\leftarrow(1-\alpha)q_o+\alpha\,\mathrm{reward}$ with
$\mathrm{reward}=1$ iff a feasible child dominates its \emph{base parent} (the primary parent the
operator perturbs; the base vector for the DE operators), else $0$, and selection
probability $p_o=p_{\min}+(1-n_{\mathrm{op}}\,p_{\min})\,q_o/\sum_{o'}q_{o'}$ over the
$n_{\mathrm{op}}$ operators. Decision-space modes (M3) are
re-estimated every $10$ generations by silhouette-selected $k$-means on $\phi(P)$ over
$[2,5]$ clusters. For structured applications $\phi$ is the problem's permutation-invariant
feature map (for placement, the $K$ stations sorted by coordinate, so two encodings of the
same layout coincide).

\subsection{Supporting modules (M2--M7)}
\textbf{M2 Three-archive system}: a Pareto-objective archive (convergence/quality), a
decision-mode archive (mode representatives), and a structure-family archive
(case-study task), with capacities $200/50/50$. Each is updated every generation by inserting new
non-dominated solutions and truncating back to capacity by crowding distance. The final reported
set (Algorithm~\ref{alg:ears}, last line) is the non-dominated union of the population and these
archives, then truncated to $N$ by the same crowding rule, so \emph{every} algorithm is evaluated
on an $N$-sized output under one algorithm-independent rule (no method gains a coverage advantage
from a larger output set). \textbf{M3 Adaptive niching}: silhouette-selected $k$-means decision-space
modes, adaptive radius, occupancy entropy; supplies the niche-rarity term and the mating
controller. \textbf{M4 Within-/cross-mode mating}: the cross-mode probability is adapted
by mode-occupancy entropy (more cross-mode exploration when modes collapse).
\textbf{M5 Constraint-aware selection}: $\epsilon$-relaxation schedule with boundary-near
preservation (a general capability, inert on the box-constrained-only problems studied here).
\textbf{M6 Operator portfolio}:
SBX+PM, DE/rand, DE/current-to-best, Gaussian and mode-interpolation operators driven by a
probability-matching bandit on credit assignment, in the adaptive-operator-selection
tradition~\parencite{dacosta2008adaptive}. \textbf{M7 Environmental selection}:
constraint-aware non-dominated sorting followed by the hybrid key of
Eq.~\eqref{eq:hybrid}.

\paragraph{Two layers: EARS-Core and EARS-Full} We distinguish two configurations to keep the
\emph{conceptual} contribution separate from the \emph{performance} of the integrated system.
\textbf{EARS-Core} is the minimal carrier of the placement idea: non-dominated sorting, the
equivalence diversity $E$, and the within-front sparsity $S$ (no archives, cross-mode mating, or
operator portfolio). It is the configuration used to test the placement claim in
Section~\ref{sec:novelty} (the minimal-skeleton rows). \textbf{EARS-Full} adds the supporting
modules M2--M7 and is the configuration ranked against the baselines. We are explicit about where
the benchmark gain comes from (the constructive ablation, Section~\ref{sec:ablation}): the
equivalence key $E$ is the dominant driver (IGDX $0.269\!\to\!0.060$), the within-front $S$ adds a
smaller but consistent increment ($\to\!0.050$), and the operator portfolio contributes further
($\to\!0.037$). The \emph{contribution we claim} is the placement of $S$ (validated on EARS-Core),
not that $S$ alone produces the rank-1 result; the rank-1 performance is a property of EARS-Full,
in which $E$ and the portfolio are the larger contributors.

\subsection{Structure-preserving diversity}
For applications whose decision space is structured (e.g.\ station layouts), EARS-MMOEA
measures decision-space diversity in a problem-provided, permutation-invariant feature
map, so $S$ in Eq.~\eqref{eq:hybrid} tracks genuine structural difference rather than raw
encoding distance. The constraint-aware non-dominated sort is vectorized via a domination
matrix at $O(N^2 m)$ for population size $N$ and $m$ objectives; the $k$NN sparsity costs
$O(N^2 d)$ for decision dimension $d$, and the periodic $k$-means re-clustering is amortized over
$10$ generations. The dominant per-generation cost is therefore $O\!\left(N^2(m{+}d)\right)$,
i.e.\ the pairwise-distance term governs and the $k$NN sparsity dominates when $d \gg m$.
